\documentclass[../main.tex]{subfiles}

\begin{document}

\chapter{Contoh Tugas, Portofolio, dan Rubrik Refleksi}

Lampiran ini memberi contoh instrumen yang dapat dipakai dosen pengampu
sekaligus contoh format bagi mahasiswa.

\section{Contoh Tugas}

\begin{enumerate}[label=\textbf{Tugas \arabic*.}, leftmargin=*]
    \item \textbf{Kriptografi klasik (Sub-CPMK 1.1).} Enkripsikan plainteks
          \code{SERANGPAGI} dengan cipher Vigen\`ere berkunci \code{KUNCI},
          kemudian dekripsikan kembali. Tunjukkan tabel konversi huruf--angka
          dan setiap langkah substitusi. Jelaskan mengapa cipher ini lebih kuat
          daripada Caesar.
    \item \textbf{Kriptografi asimetris (Sub-CPMK 1.3).} Dengan $p = 11$ dan
          $q = 13$, bangkitkan sepasang kunci RSA. Enkripsikan pesan
          $m = 7$, lalu dekripsikan cipherteksnya. Tunjukkan perhitungan
          modulus, fungsi Euler, dan invers modulo beserta pemeriksaannya.
    \item \textbf{Integritas dan otentikasi (Sub-CPMK 1.4).} Jelaskan mengapa
          fungsi hash saja tidak cukup menjamin otentikasi pesan, sehingga
          diperlukan MAC atau tanda tangan digital. Berikan satu contoh
          serangan bila hash dipakai tanpa kunci.
\end{enumerate}

\section{Panduan Portofolio}

Portofolio adalah kumpulan bukti belajar satu semester. Susunannya sebagai
berikut.

\begin{enumerate}[label=\arabic*.]
    \item \textbf{Halaman sampul}---identitas dan judul portofolio.
    \item \textbf{Daftar isi}---peta bukti yang dikumpulkan.
    \item \textbf{Jawaban latihan}---pengerjaan \textit{Latihan dan Refleksi}
          pada setiap bab, disusun menurut nomor bab.
    \item \textbf{Laporan praktikum}---hasil \textit{Aktivitas Rekomendasi}
          beserta hasil verifikasi.
    \item \textbf{Laporan proyek}---studi kasus penerapan kriptografi.
    \item \textbf{Lembar refleksi}---refleksi ketercapaian kompetensi tiap bab
          (lihat rubrik di bawah).
\end{enumerate}

\section{Rubrik Refleksi}

\begin{table}[H]
\centering
\caption{Rubrik penilaian lembar refleksi}
\label{tab:rubrik-refleksi}
\footnotesize
\begin{tabularx}{\linewidth}{|>{\raggedright\arraybackslash}X|>{\centering\arraybackslash}p{1.6cm}|>{\raggedright\arraybackslash}X|}
\hline
\textbf{Kriteria} & \textbf{Bobot} & \textbf{Indikator} \\
\hline
Kedalaman refleksi & 30\% & Menyebut apa yang sudah dan belum dipahami secara spesifik. \\
\hline
Keterkaitan dengan Sub-CPMK & 25\% & Menghubungkan refleksi dengan Sub-CPMK bab terkait. \\
\hline
Rencana perbaikan & 25\% & Menyusun langkah lanjutan yang konkret dan dapat diukur. \\
\hline
Kesadaran diri & 20\% & Menilai capaian sendiri secara jujur dan beralasan. \\
\hline
\end{tabularx}
\end{table}

Lembar refleksi yang lengkap memuat tiga bagian: apa yang dipelajari, bagian
yang masih sulit, dan rencana tindak lanjut. Ketiganya dinilai terhadap rubrik
di atas.

\end{document}
